\chapter{Introduction}
\label{ch:intro}

HPINNS is a Python package for shallow-water physics-informed neural networks.
The network represents water depth and velocity.
Training minimizes the mean square of the shallow-water residuals at a fixed grid of collocation points, together with a boundary penalty when constraints are soft.
When the selected benchmark has an exact solution, a later analysis step compares the trained field with that solution.

The package name on disk is \file{swe_pinn}.
The program you run is \file{run_pinn.py}.
Every setting that changes from one experiment to the next lives in \file{user_input.py}.

\section{What one run does}

A call to \code{python run_pinn.py} does the following.

\begin{enumerate}
\item Read \code{USER_PINN_INPUTS} from \file{user_input.py}.
\item Build the benchmark named by \code{problem_name}.
  The physical constants for that benchmark come from the matching dataclass at the bottom of the same file.
\item If a search is requested, draw trials until the completed-trial cap is reached.
  Each trial builds one network, trains it with Adam, and scores the same loss on a finer lattice.
\item Rebuild the best completed trial, or the fixed hyperparameters if the search was skipped, and train that model with Adam for \code{final_adam_epochs}.
\item Optionally continue from those weights with L-BFGS-B.
\item Write weights, a hyperparameter summary, and NumPy archives of the loss history and the predicted fields.
\end{enumerate}

Two flags split that sequence so several processes can share one study.

\begin{itemize}
\item \code{python run_pinn.py --search-only --worker-id N} runs trials only.
  The sampler seed is \(42+N\).
  The process does not train the final model.
\item \code{python run_pinn.py --finalize} loads the best completed trial already stored in the study and runs the final Adam stage, the optional L-BFGS-B stage, and the archives.
  It adds no trials.
\end{itemize}

The two flags cannot be combined.
\file{run_cores.sh} starts several search workers, waits for them, and then runs \code{--finalize} once.

\section{The three families}

Every benchmark can be trained with any of the three families.
The family is a property of the network, not of the equation.

\begin{description}
\item[\code{classical}]
Each field has its own tanh multilayer perceptron.
There is no circuit.
Width and depth set the size of each tower.

\item[\code{hybrid_quantum}]
Each field keeps that perceptron.
After the last hidden layer the tower adds a residual quantum correction,
\[
z \leftarrow z + \operatorname{decode}\!\big(\operatorname{circuit}(\operatorname{encode}(z))\big).
\]
The correction is added at full strength on every forward pass after an optional classical warmup.
A hybrid specification with zero qubits is stored and trained as classical.

\item[\code{full_quantum}]
There is no hidden perceptron and no residual skip.
Each field is
\[
\operatorname{readout}\!\big(\operatorname{circuit}(\operatorname{encode}(\boldsymbol{\xi}))\big),
\]
where \(\boldsymbol{\xi}\) is the pair of coordinates scaled onto the unit interval.
Width and depth are ignored.
The qubit count and the layer count must both be at least one.
\end{description}

Chapter~\ref{ch:families} specifies the circuit, the two embedding layouts, and the parameter counts.
Chapter~\ref{ch:examples} gives a complete fixed-architecture run for each family.
Chapter~\ref{ch:search} shows how to put all three families in one Optuna study.

\section{What this manual assumes}

You edit \file{user_input.py}, run the solver from the repository root, and read the text files and archives under \file{analysis/}.
You do not need to edit the trainer to change the family, the benchmark, the constraint mode, or the search.
A new benchmark is added by registering a problem builder, which is described briefly in Section~\ref{sec:new-problem}.
